% Part: second-order-logic
% Chapter: syntax-and-semantics
% Section: semantic-notions

\documentclass[../../../include/open-logic-section]{subfiles}

\begin{document}

\olfileid[hi]{sol}{syn}{sem}
\olsection{अर्थवैज्ञानिक धारणाएँ}

\begin{explain}
\emph{वैधता}, \emph{अनुगमन} और \emph{संतुष्ट्यता} की केंद्रीय तार्किक
धारणाएँ द्वितीय-कोटि तर्कशास्त्र के लिए उसी प्रकार परिभाषित होती हैं जैसे
प्रथम-कोटि तर्कशास्त्र के लिए; अंतर केवल यह है कि आधारभूत संतुष्टि संबंध अब
द्वितीय-कोटि !!{formula} का है। निस्संदेह, द्वितीय-कोटि !!{sentence} ऐसा
!!a{formula} है जिसमें विधेय-चर और फलन-चर सहित सभी चर आबद्ध हैं।
\end{explain}

\begin{defn}[वैधता]
वाक्य $!A$ \emph{वैध} है, $\Entails !A$, तभी और केवल तभी जब प्रत्येक
!!{structure}~$\Struct M$ के लिए $\Sat{M}{!A}$।
\end{defn}

\begin{defn}[अनुगमन]
वाक्यों का समुच्चय~$\Gamma$, वाक्य~$!A$ को \emph{अनुगमित करता है}, $\Gamma
\Entails !A$, तभी और केवल तभी जब $\Sat{M}{\Gamma}$ वाली प्रत्येक
!!{structure}~$\Struct M$ के लिए $\Sat{M}{!A}$।
\end{defn}

\begin{defn}[संतुष्ट्यता]
वाक्यों का समुच्चय~$\Gamma$ \emph{संतुष्ट्य} है यदि किसी
!!{structure}~$\Struct M$ के लिए $\Sat{M}{\Gamma}$। यदि $\Gamma$ संतुष्ट्य
नहीं है तो वह \emph{असंतुष्ट्य} कहलाता है।
\end{defn}

\end{document}
